\section{Conclusion}\label{sec:conclusion}

We formulated one step of automated physical model building as a retrieval problem: given a
description of the desired model and its input and output variables, present, from a database of
\NEq{} equations drawn mostly from the literature, the 20 equations most likely to serve in the
model.
Because an equation that contains a variable other than the inputs and outputs is useful only
together with an equation that determines that variable, the proposed two-stage method, \ours{} (two-stage retrieval with attribute-conditioned evaluation), scores each
candidate by its relation to other equations as well as by its relevance to the query.

On \NCases{} cases split by source, the method raised Recall@20 over a classical retrieval baseline
from \RecABase{} to \RecAProp{} and MAP from \MapABase{} to \MapAProp, in all 10 splits, and its
lead grew with the numbers of ground-truth equations and input variables. Ablation and permutation
importance agree on the cause: the features that read the variables shared by the query and the
equations decide most of the accuracy, the set-aware features raise the ground-truth equations
within the top 20, and the features that read topical closeness add little. The method missed
ground-truth equations only in cases for which more than 20 equations of the database hold only
query variables. Two LLMs that read the
whole database without training reached a Recall@20 of \RecAOpus{} and \RecAHaiku; \ours{} beat both in every split. The larger LLM beat the classical baseline and found ground-truth
equations that share few variables with the query, which \ours{} tends to miss; it cost
\CostPerCallOpus{} US dollars per query, and its answers varied between runs.

Three directions follow. First, the two approaches should be combined, for example by letting an
LLM rerank or extend the candidates of \ours{}. Second, the method should judge whether
two variables denote the same quantity from their definitions, not from their symbols, so that it
serves queries written in another notation. Third, the features should tell apart equations that
hold the same variables, for example by encoding which variable an equation determines, and the
benchmark should add queries that omit the intermediate variables of a model. The method should
then be joined to the combination method of \citet{CombineEq}, so that the presented equations can
be assembled into a model, solved, and compared with plant data.
